\documentclass[a4paper]{article}
\usepackage{amsmath}
\usepackage{graphicx}
\usepackage{indentfirst}
\usepackage{listings}
\usepackage{matlab-prettifier}

\author{Sergiusz Warga}
\title{Partial Differential Equations with Applications in Physics and Industry}

\begin{document}

\maketitle

\section{Visualizing characteristic Families and Decay --- Method of characteristics}

\subsection{Task A (Constant Advection)}

For $u_t + 10 u_x = 0$, write a MATLAB script to plot the parallel
characteristic lines $x(t) = x_0 + 10 t$ on the x --- t plane.
Use \texttt{plot3} or \texttt{surf} to show the initial condition
$u(x,0)=\frac{1}{1 + x^2}$ sliding along these paths without changing shape.

\section{The Riemann Problem and Shock Trajectories --- Shock Waves}

\subsection{Task A}

For $u_t + u^2 u_x = 0$ with the given step down from 2 to 1, calculate the shock speed $s$ manually.

\begin{equation}
  u(x, 0) = 
    \begin{cases}
    2 \quad \text{for} \quad x \le 0, \\
    1 \quad \text{for} \quad x > 0
    \end{cases}
\end{equation}

We can represent our Partial Differential Equation as:
\begin{align}
    u_t + u^2 u_x = 0 \\
    u_t + f(u)_x = 0 \\
    u_t + \frac{\delta f(u)}{\delta x} = 0
\end{align}

Now we find our $f(u)$.

\begin{equation}
  \frac{\delta}{\delta x}\frac{u^3}{3} = u^2 \frac{\delta u}{\delta x}
\end{equation}

\begin{equation}
  f(u) = \frac{u^3}{3}
\end{equation}


To calculate the shock speed $s$ we use Rankine-Hugoniot jump condition:
% It comes from the conservation law.
\begin{equation}
  s = \frac{f(u_R) - f(u_L)}{u_R - u_L}
\end{equation}

In our case $u_L = 2$ and $u_R = 1$.

\begin{equation}
  s = \frac{f(1) - f(2)}{1 - 2} = 2\frac{1}{3}
\end{equation}

\subsection{Task B}

Write a MATLAB script to plot the x---t plane.
Plot the characteristic lines emerging from $x < 0$ (speed 4) and $x > 0$ (speed 1).
Overlay a thick red line representing the analytical shock trajectory $x_s(t) = st$.

\begin{figure}
  \includegraphics[width=\textwidth]{figures/ex_4_b_surface.png}
  \caption{The x---t plane surface plot representing $u(x,t)$.
  The vertical lines are the numerical artifact resulting from number of points in vectors $t$ and $v$}
\end{figure}

\begin{figure}
  \includegraphics[width=\textwidth]{figures/ex_4_b_xt_plane.pdf}
\end{figure}

\lstinputlisting[style=Matlab-editor]{ex_4_b.m}

\end{document}
